Chuẩn hóa là quá trình tổ chức dữ liệu nhằm giảm dư thừa, hạn chế các bất thường khi thêm, sửa hoặc xóa dữ liệu và nâng cao tính nhất quán của cơ sở dữ liệu.

Trong đề tài này, quá trình chuẩn hóa được phân tích lần lượt qua ba mức chính: Chuẩn 1NF, Chuẩn 2NF và Chuẩn 3NF.

\section{Chuẩn thứ nhất -- 1NF}

Một quan hệ đạt chuẩn 1NF khi:

\begin{itemize}
    \item Mỗi ô dữ liệu chỉ chứa một giá trị nguyên tử.
    \item Không tồn tại nhóm thuộc tính lặp.
    \item Mỗi dòng được xác định duy nhất bởi khóa chính.
\end{itemize}

Ví dụ, bảng Customer có cấu trúc:

Customer(customer\_id, full\_name, phone, email, address)

Mỗi thuộc tính như full\_name, phone, email và address chỉ chứa một giá trị tại mỗi bản ghi. Không có thuộc tính dạng danh sách hoặc nhóm giá trị lặp.

Tương tự, các bảng Pet, Veterinarian, Staff, Admin, Booking, Invoice, MedicalRecord và Room đều được thiết kế với các thuộc tính nguyên tử.

\begin{table}[htbp]
    \centering
    \caption{Kiểm tra chuẩn 1NF}
    \label{tab:1nf}
    \begin{tabularx}{\textwidth}{|p{3cm}|X|p{3cm}|}
        \hline
        \textbf{Quan hệ} & \textbf{Kiểm tra} & \textbf{Kết quả} \\ \hline
        Customer & Thuộc tính nguyên tử, không có nhóm lặp & Đạt 1NF \\ \hline
        Pet & Thuộc tính nguyên tử, không có nhóm lặp & Đạt 1NF \\ \hline
        Veterinarian & Thuộc tính nguyên tử, không có nhóm lặp & Đạt 1NF \\ \hline
        Staff & Thuộc tính nguyên tử, không có nhóm lặp & Đạt 1NF \\ \hline
        Admin & Thuộc tính nguyên tử, không có nhóm lặp & Đạt 1NF \\ \hline
        Booking & Thuộc tính nguyên tử, không có nhóm lặp & Đạt 1NF \\ \hline
        Invoice & Thuộc tính nguyên tử, không có nhóm lặp & Đạt 1NF \\ \hline
        MedicalRecord & Thuộc tính nguyên tử, không có nhóm lặp & Đạt 1NF \\ \hline
        Room & Thuộc tính nguyên tử, không có nhóm lặp & Đạt 1NF \\ \hline
    \end{tabularx}
\end{table}

\section{Chuẩn thứ hai -- 2NF}

Một quan hệ đạt chuẩn 2NF khi:

\begin{enumerate}
    \item Đã đạt 1NF.
    \item Mọi thuộc tính không khóa phải phụ thuộc đầy đủ vào toàn bộ khóa chính.
\end{enumerate}

Trong cơ sở dữ liệu hiện tại, tất cả các bảng đều sử dụng khóa chính đơn, tức khóa chính chỉ gồm một thuộc tính.

Ví dụ:

\[
Customer:
customer\_id \rightarrow name, phone, email, address
\]

Do khóa chính chỉ gồm customer\_id nên không thể tồn tại phụ thuộc bộ phận.

Tương tự:

\[
Pet:
pet\_id \rightarrow name, species, breed, age, customer\_id
\]

\[
Booking:
booking\_id \rightarrow date\_time, status, payment,
customer\_id, pet\_id, vet\_id, staff\_id
\]

Do đó, các quan hệ đều không có phụ thuộc bộ phận và đạt chuẩn 2NF.

\begin{table}[htbp]
    \centering
    \caption{Kiểm tra chuẩn 2NF}
    \label{tab:2nf}
    \begin{tabularx}{\textwidth}{|X|X|X|}
        \hline
        \textbf{Quan hệ} & \textbf{Đặc điểm khóa} & \textbf{Kết quả} \\ \hline
        Customer & Khóa đơn & Đạt 2NF \\ \hline
        Pet & Khóa đơn & Đạt 2NF \\ \hline
        Veterinarian & Khóa đơn & Đạt 2NF \\ \hline
        Staff & Khóa đơn & Đạt 2NF \\ \hline
        Admin & Khóa đơn & Đạt 2NF \\ \hline
        Booking & Khóa đơn & Đạt 2NF \\ \hline
        Invoice & Khóa đơn & Đạt 2NF \\ \hline
        MedicalRecord & Khóa đơn & Đạt 2NF \\ \hline
        Room & Khóa đơn & Đạt 2NF \\ \hline
    \end{tabularx}
\end{table}

\section{Chuẩn thứ ba -- 3NF}

Một quan hệ đạt chuẩn 3NF khi:

\begin{enumerate}
    \item Đã đạt 2NF.
    \item Không tồn tại phụ thuộc bắc cầu từ khóa chính đến thuộc tính không khóa thông qua một thuộc tính không khóa khác.
\end{enumerate}

Ví dụ với bảng Customer:

\[
customer\_id \rightarrow full\_name, phone, email, address
\]

Các thuộc tính full\_name, phone, email và address phụ thuộc trực tiếp vào customer\_id. Không có thuộc tính không khóa nào xác định một thuộc tính không khóa khác trong bảng.

Đối với Booking:

\[
booking\_id \rightarrow date\_time, status, payment,
customer\_id, pet\_id, vet\_id, staff\_id
\]

Thông tin tên khách hàng không được lưu trực tiếp trong Booking. Thay vào đó, Booking chỉ lưu customer\_id và thông tin khách hàng được truy xuất từ bảng Customer.

Tương tự, Booking chỉ lưu vet\_id thay vì lưu tên và chuyên môn của bác sĩ. Điều này giúp tránh phụ thuộc bắc cầu và dữ liệu dư thừa.

\section{Tổng hợp kết quả chuẩn hóa}

\begin{table}[htbp]
    \centering
    \caption{Kết quả chuẩn hóa cơ sở dữ liệu}
    \label{tab:normalization_summary}
    \begin{tabularx}{\textwidth}{|X|X|X|X|}
        \hline
        \textbf{Quan hệ} &
        \textbf{1NF} &
        \textbf{2NF} &
        \textbf{3NF} \\ \hline

        Customer & Đạt & Đạt & Đạt \\ \hline
        Pet & Đạt & Đạt & Đạt \\ \hline
        Veterinarian & Đạt & Đạt & Đạt \\ \hline
        Staff & Đạt & Đạt & Đạt \\ \hline
        Admin & Đạt & Đạt & Đạt \\ \hline
        Booking & Đạt & Đạt & Đạt \\ \hline
        Invoice & Đạt & Đạt & Đạt \\ \hline
        MedicalRecord & Đạt & Đạt & Đạt \\ \hline
        Room & Đạt & Đạt & Đạt \\ \hline
    \end{tabularx}
\end{table}

\section{Các bất thường được hạn chế}

Việc chuẩn hóa giúp hạn chế ba loại bất thường phổ biến.

\begin{table}[htbp]
    \centering
    \caption{Các bất thường trong cơ sở dữ liệu}
    \label{tab:data_anomalies}
    \begin{tabularx}{\textwidth}{|p{4cm}|X|}
        \hline
        \textbf{Loại bất thường} & \textbf{Mô tả} \\ \hline

        Khi thêm dữ liệu &
        Thông tin khách hàng có thể được thêm độc lập mà không cần tạo booking. Tương tự, thông tin bác sĩ có thể được quản lý độc lập.
        \\ \hline

        Khi cập nhật &
        Thông tin khách hàng chỉ được lưu tại bảng Customer. Khi thay đổi số điện thoại hoặc địa chỉ, chỉ cần cập nhật một bản ghi thay vì cập nhật nhiều bản ghi booking.
        \\ \hline

        Khi xóa &
        Việc xóa một booking không làm mất thông tin khách hàng hoặc thông tin bác sĩ vì các dữ liệu này được lưu ở các quan hệ riêng biệt.
        \\ \hline

    \end{tabularx}
\end{table}

\section{Kết luận chương}

Cơ sở dữ liệu được phân tích lần lượt theo ba mức chuẩn 1NF, 2NF và 3NF. Các thuộc tính đều có tính nguyên tử, các quan hệ không có phụ thuộc bộ phận và không tồn tại phụ thuộc bắc cầu không cần thiết. Thiết kế hiện tại đạt chuẩn 3NF và hạn chế đáng kể dữ liệu dư thừa cũng như các bất thường trong quá trình thao tác dữ liệu.